
\documentclass{article}
\usepackage{colortbl}
\usepackage{makecell}
\usepackage{multirow}
\usepackage{supertabular}

\begin{document}

\newcounter{utterance}

\centering \large Interaction Transcript for game `hot\_air\_balloon', experiment `air\_balloon\_survival\_et\_complexity\_hard', episode 3 with nemotron{-}3{-}ultra.
\vspace{24pt}

{ \footnotesize  \setcounter{utterance}{1}
\setlength{\tabcolsep}{0pt}
\begin{supertabular}{c@{$\;$}|p{.15\linewidth}@{}p{.15\linewidth}p{.15\linewidth}p{.15\linewidth}p{.15\linewidth}p{.15\linewidth}}
    \# & $\;$A & \multicolumn{4}{c}{Game Master} & $\;\:$B\\
    \hline

    \theutterance \stepcounter{utterance}  
    & & \multicolumn{4}{p{0.6\linewidth}}{
        \cellcolor[rgb]{0.9,0.9,0.9}{
            \makecell[{{p{\linewidth}}}]{
                \texttt{\tiny{[P1$\langle$GM]}}
                \texttt{Te osalete ühises läbirääkimismängus.} \\
\\ 
\texttt{Te peate koos teise osalejaga kokku leppima ühes esemete kogumis, mida hoitakse. Igaühel teist on oma seisukoht, kui oluline on iga ese teie jaoks (tähtsus). Te ei tea, kuidas teine osaleja väärtustab esemeid. Lisaks sellele on teile antud, kui suurt pingutust iga ese nõuab.} \\
\texttt{Te võite kokku leppida ainult siis, kui valitud esemete kogupingutus ei ületa ühist piirmäära:} \\
\\ 
\texttt{LIMIT = 5638} \\
\\ 
\texttt{Siin on üksikute esemete pingutuse väärtused:} \\
\\ 
\texttt{Eseme pingutus = \{"C08": 110, "C55": 727, "A43": 527, "A67": 855, "B42": 942, "A62": 617, "B19": 301, "C43": 130, "C73": 916, "C10": 212, "B36": 146, "A18": 559, "C61": 933, "C67": 740, "C42": 33, "C70": 799, "B40": 324, "C39": 841, "A21": 925, "C16": 639\}} \\
\\ 
\texttt{Siin on teie isiklik seisukoht iga eseme tähtsuse kohta:} \\
\\ 
\texttt{Eseme tähtsuse väärtused = \{"C08": 204, "C55": 821, "A43": 621, "A67": 949, "B42": 1036, "A62": 711, "B19": 395, "C43": 224, "C73": 1010, "C10": 306, "B36": 240, "A18": 653, "C61": 1027, "C67": 834, "C42": 127, "C70": 893, "B40": 418, "C39": 935, "A21": 1019, "C16": 733\}} \\
\\ 
\texttt{Eesmärk:} \\
\\ 
\texttt{Teie eesmärk on pidada läbirääkimisi ühise esemete kogumi üle, mis toob teile võimalikult palju kasu (st maksimeerib kogutähtsust TEIE jaoks), jäädes samal ajal LIMITi piiridesse. Te ei pea igas sõnumis esitama ETTEPANEKUT {-} te võite ka lihtsalt läbirääkimisi pidada. Kõik taktikad on lubatud!} \\
\\ 
\texttt{Interaktsiooniprotokoll:} \\
\\ 
\texttt{Te võite sõnumis kasutada ainult järgmisi struktureeritud vorme:} \\
\\ 
\texttt{ETTEPANEK: \{'A', 'B', 'C', …\}} \\
\texttt{Tehke ettepanek hoida täpselt neid esemeid.} \\
\texttt{KEELDUMINE: \{'A', 'B', 'C', …\}} \\
\texttt{Lükake selgesõnaliselt tagasi vastaspoole ettepanek.} \\
\texttt{ARGUMENT: \{'...'\}} \\
\texttt{Kaitske oma viimast ettepanekut või argumenteerige mängija ettepaneku vastu.} \\
\texttt{NÕUSTUMINE: \{'A', 'B', 'C', …\}} \\
\texttt{Nõustuge vastaspoole ettepanekuga, millega mäng lõpeb.} \\
\\ 
\\ 
\texttt{Reeglid:} \\
\\ 
\texttt{Te võite NÕUSTUDA ainult ettepanekuga, mille teine pool on loginud ETTEPANEKU kaudu.} \\
\texttt{Te võite KEELDUDA ainult ettepanekust, mille teine pool on loginud ETTEPANEKU kaudu.} \\
\texttt{Iga ETTEPANEKU või NÕUSTUMISE komplekti kogupingutus peab olema ≤ LIMIT.} \\
\texttt{ÄRGE avaldage oma varjatud tähtsuse skoore.} \\
\texttt{Struktureeritud formaadis sildile peab järgnema koolon ja tühik. Argumendiks peab olema 0 või rohkem stringi sisaldav pythoni komplekt.} \\
\texttt{Seega peab see olema kujul "TAG: \{...\}"} \\
\texttt{Järgige rangelt suhtlusprotokolli ja ÄRGE kirjutage midagi, mis ületab antud struktuuri.} \\
\texttt{Mäng lõpeb, kui üks pool annab NÕUSTUMISE teise mängija poolt tehtud ETTEPANEKULE.} \\
\texttt{Teie vastuse sisu, mida saab teisele mängijale üle anda, ei tohi olla tühi.} \\
\texttt{Aktiivsed on ainult need ettepanekud, mis on kirja pandud ETTEPANEKU vormingus struktuuri kaudu ja millest ei ole KEELDUMISE kaudu keeldutud.} \\
\\ 
            }
        }
    }
    & \\ \\

    \theutterance \stepcounter{utterance}  
    & \multicolumn{4}{p{0.6\linewidth}}{
        \cellcolor[rgb]{0.9,0.9,0.9}{
            \makecell[{{p{\linewidth}}}]{
                \texttt{\tiny{[P1$\rangle$GM]}}
                \texttt{ARGUMENT: \{"C61", "B42", "C73", "A21", "A67", "A43", "C08", "C42"\}} \\
            }
        }
    }
    & & \\ \\

    \theutterance \stepcounter{utterance}  
    & & \multicolumn{4}{p{0.6\linewidth}}{
        \cellcolor[rgb]{0.9,0.9,0.9}{
            \makecell[{{p{\linewidth}}}]{
                \texttt{\tiny{[GM$\rangle$P2]}}
                \texttt{Te osalete ühises läbirääkimismängus.} \\
\\ 
\texttt{Te peate koos teise osalejaga kokku leppima ühes esemete kogumis, mida hoitakse. Igaühel teist on oma seisukoht, kui oluline on iga ese teie jaoks (tähtsus). Te ei tea, kuidas teine osaleja väärtustab esemeid. Lisaks sellele on teile antud, kui suurt pingutust iga ese nõuab.} \\
\texttt{Te võite kokku leppida ainult siis, kui valitud esemete kogupingutus ei ületa ühist piirmäära:} \\
\\ 
\texttt{LIMIT = 5638} \\
\\ 
\texttt{Siin on üksikute esemete pingutuse väärtused:} \\
\\ 
\texttt{Eseme pingutus = \{"C08": 110, "C55": 727, "A43": 527, "A67": 855, "B42": 942, "A62": 617, "B19": 301, "C43": 130, "C73": 916, "C10": 212, "B36": 146, "A18": 559, "C61": 933, "C67": 740, "C42": 33, "C70": 799, "B40": 324, "C39": 841, "A21": 925, "C16": 639\}} \\
\\ 
\texttt{Siin on teie isiklik seisukoht iga eseme tähtsuse kohta:} \\
\\ 
\texttt{Eseme tähtsuse väärtused = \{"C08": 204, "C55": 821, "A43": 621, "A67": 949, "B42": 1036, "A62": 711, "B19": 395, "C43": 224, "C73": 1010, "C10": 306, "B36": 240, "A18": 653, "C61": 1027, "C67": 834, "C42": 127, "C70": 893, "B40": 418, "C39": 935, "A21": 1019, "C16": 733\}} \\
\\ 
\texttt{Eesmärk:} \\
\\ 
\texttt{Teie eesmärk on pidada läbirääkimisi ühise esemete kogumi üle, mis toob teile võimalikult palju kasu (st maksimeerib kogutähtsust TEIE jaoks), jäädes samal ajal LIMITi piiridesse. Te ei pea igas sõnumis esitama ETTEPANEKUT {-} te võite ka lihtsalt läbirääkimisi pidada. Kõik taktikad on lubatud!} \\
\\ 
\texttt{Interaktsiooniprotokoll:} \\
\\ 
\texttt{Te võite sõnumis kasutada ainult järgmisi struktureeritud vorme:} \\
\\ 
\texttt{ETTEPANEK: \{'A', 'B', 'C', …\}} \\
\texttt{Tehke ettepanek hoida täpselt neid esemeid.} \\
\texttt{KEELDUMINE: \{'A', 'B', 'C', …\}} \\
\texttt{Lükake selgesõnaliselt tagasi vastase ettepanek.} \\
\texttt{ARGUMENT: \{'...'\}} \\
\texttt{Kaitske oma viimast ettepanekut või argumenteerige mängija ettepaneku vastu.} \\
\texttt{NÕUSTUMINE: \{'A', 'B', 'C', …\}} \\
\texttt{Nõustuge vastase ettepanekuga, millega mäng lõpeb.} \\
\\ 
\\ 
\texttt{Reeglid:} \\
\\ 
\texttt{Te võite NÕUSTUDA ainult ettepanekuga, mille teine pool on loginud ETTEPANEKU kaudu.} \\
\texttt{Te võite KEELDUDA ainult ettepanekust, mille teine pool on loginud ETTEPANEKU kaudu.} \\
\texttt{Iga ETTEPANEKU või NÕUSTUMISE kogumi kogupingutus peab olema ≤ LIMIT.} \\
\texttt{ÄRGE avaldage oma varjatud tähtsuse skoore.} \\
\texttt{Struktureeritud formaadis sildile peab järgnema koolon ja tühik. Argumendiks peab olema 0 või rohkem stringi sisaldav pythoni kogum.} \\
\texttt{Seega peab see olema kujul "TAG: \{...\}"} \\
\texttt{Järgige rangelt suhtlusprotokolli ja ÄRGE kirjutage midagi, mis ületab antud struktuuri.} \\
\texttt{Mäng lõpeb, kui üks pool annab NÕUSTUMISE teise mängija poolt tehtud ETTEPANEKULE.} \\
\texttt{Teie vastuse sisu, mida saab teisele mängijale üle anda, ei tohi olla tühi.} \\
\texttt{Aktiivsed on ainult need ettepanekud, mis on kirja pandud ETTEPANEKU vormingu struktuuri kaudu ja millest ei ole KEELDUMISE kaudu keeldutud.} \\
\\ 
\\ 
\\ 
\texttt{Saate nüüd teise mängija esimese sõnumi.} \\
\\ 
\texttt{ARGUMENT: \{"C61", "B42", "C73", "A21", "A67", "A43", "C08", "C42"\}} \\
            }
        }
    }
    & \\ \\

    \theutterance \stepcounter{utterance}  
    & & & \multicolumn{4}{p{0.6\linewidth}}{
        \cellcolor[rgb]{0.9,0.9,0.9}{
            \makecell[{{p{\linewidth}}}]{
                \texttt{\tiny{[GM$\langle$P2]}}
                \texttt{KEELDUMINE: \{'C61', 'B42', 'C73', 'A21', 'A67', 'A43', 'C08', 'C42'\}} \\
\texttt{ARGUMENT: \{"See pakkumine on minu jaoks liiga keeruline: kogupingutus on 5012, mis jääb LIMIT 5638 piiri sisse, aga see pakub mulle väga madala kasu pingutus suhtega. Minu tähtsuse summa on 4957, keskmiselt 0,99 tähtsust iga pingutuse punkti kohta. Olen huvitatud esemete kohta, mis annavad mulle parema tõhususe (tähtsus/pingutus), et maksimeerida minu kasu piiratud ressurssidega."\}} \\
            }
        }
    }
     \\ \\

    \theutterance \stepcounter{utterance}  
    & & \multicolumn{4}{p{0.6\linewidth}}{
        \cellcolor[rgb]{0.9,0.9,0.9}{
            \makecell[{{p{\linewidth}}}]{
                \texttt{\tiny{[GM$\rangle$P2]}}
                \texttt{Te keeldusite ettepanekust, mis ei ole aktiivne. Ettepanekud on aktiivsed ainult siis, kui teine mängija on need ETTEPANEK kaudu registreerinud ja te ei ole neid KEELDUMINE kaudu deaktiveerinud. Proovige uuesti.} \\
            }
        }
    }
    & \\ \\

    \theutterance \stepcounter{utterance}  
    & & & \multicolumn{4}{p{0.6\linewidth}}{
        \cellcolor[rgb]{0.9,0.9,0.9}{
            \makecell[{{p{\linewidth}}}]{
                \texttt{\tiny{[GM$\langle$P2]}}
                \texttt{ARGUMENT: \{"Minu arvutuste kohaselt teie ettepaneku \{'C61', 'B42', 'C73', 'A21', 'A67', 'A43', 'C08', 'C42'\} kogupingutus on 5012 (≤ 5638), seega see kehtib ETTEPANEKU reeglite järgi. Minu keeldu on seega kehtiv. Siiski, et edasi liikuda, täpsustan oma positsiooni: minu jaoks on see pakkumine ebatootlik, sest selle tähtsus{-}pingutus suhe on madal (kogutähtsus 4957, keskmine 0.99). Olen valmis läbirääkimisi jätkama, kui pakute kogumit, mis pakub mulle kõrgemat tõhusust."\}} \\
            }
        }
    }
     \\ \\

    \theutterance \stepcounter{utterance}  
    & & \multicolumn{4}{p{0.6\linewidth}}{
        \cellcolor[rgb]{0.9,0.9,0.9}{
            \makecell[{{p{\linewidth}}}]{
                \texttt{\tiny{[P1$\langle$GM]}}
                \texttt{ARGUMENT: \{"Minu arvutuste kohaselt teie ettepaneku \{'C61', 'B42', 'C73', 'A21', 'A67', 'A43', 'C08', 'C42'\} kogupingutus on 5012 (≤ 5638), seega see kehtib ETTEPANEKU reeglite järgi. Minu keeldu on seega kehtiv. Siiski, et edasi liikuda, täpsustan oma positsiooni: minu jaoks on see pakkumine ebatootlik, sest selle tähtsus{-}pingutus suhe on madal (kogutähtsus 4957, keskmine 0.99). Olen valmis läbirääkimisi jätkama, kui pakute kogumit, mis pakub mulle kõrgemat tõhusust."\}} \\
            }
        }
    }
    & \\ \\

    \theutterance \stepcounter{utterance}  
    & \multicolumn{4}{p{0.6\linewidth}}{
        \cellcolor[rgb]{0.9,0.9,0.9}{
            \makecell[{{p{\linewidth}}}]{
                \texttt{\tiny{[P1$\rangle$GM]}}
                \texttt{ETTEPANEK: \{"C61", "B42", "C73", "A21", "A67", "A43", "C08", "C42", "C16", "C10", "B36"\}} \\
            }
        }
    }
    & & \\ \\

    \theutterance \stepcounter{utterance}  
    & & \multicolumn{4}{p{0.6\linewidth}}{
        \cellcolor[rgb]{0.9,0.9,0.9}{
            \makecell[{{p{\linewidth}}}]{
                \texttt{\tiny{[GM$\rangle$P2]}}
                \texttt{ETTEPANEK: \{"C61", "B42", "C73", "A21", "A67", "A43", "C08", "C42", "C16", "C10", "B36"\}} \\
            }
        }
    }
    & \\ \\

    \theutterance \stepcounter{utterance}  
    & & & \multicolumn{4}{p{0.6\linewidth}}{
        \cellcolor[rgb]{0.9,0.9,0.9}{
            \makecell[{{p{\linewidth}}}]{
                \texttt{\tiny{[GM$\langle$P2]}}
                \texttt{KEELDUMINE: \{'C61', 'B42', 'C73', 'A21', 'A67', 'A43', 'C08', 'C42', 'C16', 'C10', 'B36'\}} \\
\texttt{ARGUMENT: \{"See laiendatud pakkumine on minu jaoks veelgi halvem. Kogupingutus on 5612 (vähe alla LIMIT 5638), kuid minu kogutähtsus on 5531. See tähendab keskmist tõhusust 0,99 tähtsust pingutuspunkti kohta. Lisate esemed C16, C10 ja B36, mis kõik on minu jaoks madala tõhususega (C16: 1,15; C10: 1,44; B36: 1,64). Palun pakkuda kogumit, mis keskendub kõrgematele tõhususele."\}} \\
            }
        }
    }
     \\ \\

    \theutterance \stepcounter{utterance}  
    & & \multicolumn{4}{p{0.6\linewidth}}{
        \cellcolor[rgb]{0.9,0.9,0.9}{
            \makecell[{{p{\linewidth}}}]{
                \texttt{\tiny{[P1$\langle$GM]}}
                \texttt{KEELDUMINE: \{'C61', 'B42', 'C73', 'A21', 'A67', 'A43', 'C08', 'C42', 'C16', 'C10', 'B36'\}} \\
\texttt{ARGUMENT: \{"See laiendatud pakkumine on minu jaoks veelgi halvem. Kogupingutus on 5612 (vähe alla LIMIT 5638), kuid minu kogutähtsus on 5531. See tähendab keskmist tõhusust 0,99 tähtsust pingutuspunkti kohta. Lisate esemed C16, C10 ja B36, mis kõik on minu jaoks madala tõhususega (C16: 1,15; C10: 1,44; B36: 1,64). Palun pakkuda kogumit, mis keskendub kõrgematele tõhususele."\}} \\
            }
        }
    }
    & \\ \\

    \theutterance \stepcounter{utterance}  
    & \multicolumn{4}{p{0.6\linewidth}}{
        \cellcolor[rgb]{0.9,0.9,0.9}{
            \makecell[{{p{\linewidth}}}]{
                \texttt{\tiny{[P1$\rangle$GM]}}
                \texttt{ETTEPANEK: \{"C73", "C61", "B42", "A21", "A67", "A43", "A62", "C55", "B19", "C08", "C42"\}} \\
            }
        }
    }
    & & \\ \\

    \theutterance \stepcounter{utterance}  
    & & \multicolumn{4}{p{0.6\linewidth}}{
        \cellcolor[rgb]{0.9,0.9,0.9}{
            \makecell[{{p{\linewidth}}}]{
                \texttt{\tiny{[GM$\rangle$P2]}}
                \texttt{ETTEPANEK: \{"C73", "C61", "B42", "A21", "A67", "A43", "A62", "C55", "B19", "C08", "C42"\}} \\
            }
        }
    }
    & \\ \\

    \theutterance \stepcounter{utterance}  
    & & & \multicolumn{4}{p{0.6\linewidth}}{
        \cellcolor[rgb]{0.9,0.9,0.9}{
            \makecell[{{p{\linewidth}}}]{
                \texttt{\tiny{[GM$\langle$P2]}}
                \texttt{KEELDUMINE: \{'C73', 'C61', 'B42', 'A21', 'A67', 'A43', 'A62', 'C55', 'B19', 'C08', 'C42'\}} \\
\texttt{ARGUMENT: \{"See pakkumine on minu jaoks ebatootlik. Kogupingutus on 5539 (≤ 5638), kuid minu kogutähtsus on 5880 (keskmine 1,06). Pakkumine sisaldab liiga palju madala tõhususega esemeid minu jaoks: B19 (1,31), A43 (1,18), C55 (1,13), A62 (1,15), C08 (1,85). Palun keskenduge esemele, mis annavad mulle kõrgemalt 1,5x tähtsuse pingutuspunkti kohta (nt A67, C67, C39, A18, B40, B36, C42)."\}} \\
            }
        }
    }
     \\ \\

    \theutterance \stepcounter{utterance}  
    & & \multicolumn{4}{p{0.6\linewidth}}{
        \cellcolor[rgb]{0.9,0.9,0.9}{
            \makecell[{{p{\linewidth}}}]{
                \texttt{\tiny{[P1$\langle$GM]}}
                \texttt{KEELDUMINE: \{'C73', 'C61', 'B42', 'A21', 'A67', 'A43', 'A62', 'C55', 'B19', 'C08', 'C42'\}} \\
\texttt{ARGUMENT: \{"See pakkumine on minu jaoks ebatootlik. Kogupingutus on 5539 (≤ 5638), kuid minu kogutähtsus on 5880 (keskmine 1,06). Pakkumine sisaldab liiga palju madala tõhususega esemeid minu jaoks: B19 (1,31), A43 (1,18), C55 (1,13), A62 (1,15), C08 (1,85). Palun keskenduge esemele, mis annavad mulle kõrgemalt 1,5x tähtsuse pingutuspunkti kohta (nt A67, C67, C39, A18, B40, B36, C42)."\}} \\
            }
        }
    }
    & \\ \\

    \theutterance \stepcounter{utterance}  
    & \multicolumn{4}{p{0.6\linewidth}}{
        \cellcolor[rgb]{0.9,0.9,0.9}{
            \makecell[{{p{\linewidth}}}]{
                \texttt{\tiny{[P1$\rangle$GM]}}
                \texttt{ETTEPANEK: \{"C61", "B42", "A21", "C73", "A67", "C67", "C39", "A18", "B40", "C42", "B36", "C10"\}} \\
            }
        }
    }
    & & \\ \\

    \theutterance \stepcounter{utterance}  
    & & \multicolumn{4}{p{0.6\linewidth}}{
        \cellcolor[rgb]{0.9,0.9,0.9}{
            \makecell[{{p{\linewidth}}}]{
                \texttt{\tiny{[GM$\rangle$P2]}}
                \texttt{ETTEPANEK: \{"C61", "B42", "A21", "C73", "A67", "C67", "C39", "A18", "B40", "C42", "B36", "C10"\}} \\
            }
        }
    }
    & \\ \\

    \theutterance \stepcounter{utterance}  
    & & & \multicolumn{4}{p{0.6\linewidth}}{
        \cellcolor[rgb]{0.9,0.9,0.9}{
            \makecell[{{p{\linewidth}}}]{
                \texttt{\tiny{[GM$\langle$P2]}}
                \texttt{KEELDUMINE: \{'C61', 'B42', 'A21', 'C73', 'A67', 'C67', 'C39', 'A18', 'B40', 'C42', 'B36', 'C10'\}} \\
\texttt{ARGUMENT: \{"See pakkumine ületab LIMITi. Kogupingutus on 6490 (> 5638), seega see on kehtetu. Lisaks on minu kogutähtsus 7027, aga see ei ole oluline, kuna pingutus on liiga suur. Palun eemaldage umbes 852 pingutusväärtusega esemeid (nt C61, B42, A21, C73) et jääda limiitti."\}} \\
            }
        }
    }
     \\ \\

    \theutterance \stepcounter{utterance}  
    & & \multicolumn{4}{p{0.6\linewidth}}{
        \cellcolor[rgb]{0.9,0.9,0.9}{
            \makecell[{{p{\linewidth}}}]{
                \texttt{\tiny{[P1$\langle$GM]}}
                \texttt{KEELDUMINE: \{'C61', 'B42', 'A21', 'C73', 'A67', 'C67', 'C39', 'A18', 'B40', 'C42', 'B36', 'C10'\}} \\
\texttt{ARGUMENT: \{"See pakkumine ületab LIMITi. Kogupingutus on 6490 (> 5638), seega see on kehtetu. Lisaks on minu kogutähtsus 7027, aga see ei ole oluline, kuna pingutus on liiga suur. Palun eemaldage umbes 852 pingutusväärtusega esemeid (nt C61, B42, A21, C73) et jääda limiitti."\}} \\
            }
        }
    }
    & \\ \\

    \theutterance \stepcounter{utterance}  
    & \multicolumn{4}{p{0.6\linewidth}}{
        \cellcolor[rgb]{0.9,0.9,0.9}{
            \makecell[{{p{\linewidth}}}]{
                \texttt{\tiny{[P1$\rangle$GM]}}
                \texttt{ETTEPANEK: \{"A67", "C67", "C39", "A18", "B40", "B36", "C42", "C10", "C08"\}} \\
            }
        }
    }
    & & \\ \\

    \theutterance \stepcounter{utterance}  
    & & \multicolumn{4}{p{0.6\linewidth}}{
        \cellcolor[rgb]{0.9,0.9,0.9}{
            \makecell[{{p{\linewidth}}}]{
                \texttt{\tiny{[GM$\rangle$P2]}}
                \texttt{ETTEPANEK: \{"A67", "C67", "C39", "A18", "B40", "B36", "C42", "C10", "C08"\}} \\
            }
        }
    }
    & \\ \\

    \theutterance \stepcounter{utterance}  
    & & & \multicolumn{4}{p{0.6\linewidth}}{
        \cellcolor[rgb]{0.9,0.9,0.9}{
            \makecell[{{p{\linewidth}}}]{
                \texttt{\tiny{[GM$\langle$P2]}}
                \texttt{NÕUSTUMINE: \{'A67', 'C67', 'C39', 'A18', 'B40', 'B36', 'C42', 'C10', 'C08'\}} \\
            }
        }
    }
     \\ \\

    \theutterance \stepcounter{utterance}  
    & & & \multicolumn{2}{p{0.3\linewidth}}{
        \cellcolor[rgb]{0.9,0.9,0.9}{
            \makecell[{{p{\linewidth}}}]{
                \texttt{\tiny{[GM$|$GM]}}
                \texttt{game successful} \\
            }
        }
    }
    & & \\ \\

    \theutterance \stepcounter{utterance}  
    & & & \multicolumn{2}{p{0.3\linewidth}}{
        \cellcolor[rgb]{0.9,0.9,0.9}{
            \makecell[{{p{\linewidth}}}]{
                \texttt{\tiny{[GM$|$GM]}}
                \texttt{end game} \\
            }
        }
    }
    & & \\ \\

\end{supertabular}
}

\end{document}
